\begin{theorem}
\label{theorem-BL-glueing}
\begin{reference}
Slight generalization of the main theorem of \cite{Beauville-Laszlo}.
\end{reference}
Let $(R \to R',f)$ be a glueing pair. The functor
$\text{Can} : \text{Mod}_R \longrightarrow \text{Glue}(R \to R', f)$
determines an equivalence of the category of $R$-modules glueable
for $(R \to R', f)$ and the category $\text{Glue}(R \to R', f)$
of glueing data.
\end{theorem}

\begin{proof}
Let $(M', M_1, \alpha_1)$ be a glueing datum. We will show that
$M = H^0((M', M_1, \alpha_1))$ is a glueable module for $(R \to R', f)$
and that $(M', M_1, \alpha_1) \cong \text{Can}(M)$.

\medskip\noindent
We first check that the map $\text{d} : M' \oplus M_1 \to (M')_f$ used in the
definition of the functor $H^0$ is surjective. Observe that
$(x, y) \in M' \oplus M_1$ maps to
$\text{d}(x, y) = x/1 - \alpha_1^{-1}(y \otimes 1)$
in $(M')_f$. If $z \in (M')_f$, then we can write
$\alpha_1(z) = \sum y_i \otimes g_i$ with $g_i \in R'$
and $y_i \in M_1$. Write $\alpha_1^{-1}(y_i \otimes 1) = y_i'/f^n$
for some $y'_i \in M'$ and $n \geq 0$ (we can pick the same $n$
for all $i$). Write $g_i = a_i + f^n b_i$ with $a_i \in R$ and
$b_i \in R'$. Then with $y = \sum a_i y_i \in M_1$ and
$x = \sum b_i y'_i \in M'$ we have $\text{d}(x, -y) = z$
as desired.

\medskip\noindent
Since $M = H^0((M', M_1, \alpha_1)) = \Ker(\text{d})$ we obtain
an exact sequence of $R$-modules
\begin{equation}
\label{equation-define-M}
0 \to M \to M' \oplus M_1 \to (M')_f \to 0.
\end{equation}
We will prove that the maps $M \to M'$ and $M \to M_1$ induce isomorphisms
$M \otimes_R R' \to M'$ and $M \otimes_R R_f \to M_1$.
This will imply that $M$ is glueable for $(R \to R', f)$ and
$\text{Can}(M) \cong (M', M_1, \alpha_1)$ as desired.

\medskip\noindent
Since $f$ is a nonzerodivisor on $M_1$, we have
$M[f^\infty] \cong M'[f^\infty]$. This yields an exact sequence
\begin{equation}
\label{equation-exact-mod-torsion}
0 \to M/M[f^\infty] \to M_1 \to (M')_f/M' \to 0.
\end{equation}
Since $R \to R_f$ is flat, we may tensor this exact sequence with $R_f$
to deduce that $M \otimes_R R_f = (M/M[f^\infty]) \otimes_R R_f \to M_1$
is an isomorphism.

\medskip\noindent
By Lemma \ref{lemma-BL3} we have
$\text{Tor}_1^R(R', \Coker(M' \to (M')_f)) = 0$.
The sequence (\ref{equation-exact-mod-torsion})
thus remains exact upon tensoring
over $R$ with $R'$. Using $\alpha_1$ and
Lemma \ref{lemma-torsion-completion}
the resulting exact sequence can be written as
\begin{equation}
\label{equation-mod-torsion-sequence}
0 \to (M/M[f^\infty]) \otimes_R R' \to
(M')_f \to (M')_f/M' \to 0
\end{equation}
This yields an isomorphism
$(M/M[f^\infty]) \otimes_R R' \cong M'/M'[f^\infty]$.
This implies that in the diagram
$$
\xymatrix{
& M[f^\infty] \otimes_R R' \ar[r] \ar[d] &
M \otimes_R R'  \ar[r] \ar[d] &
(M/M[f^\infty]) \otimes_R R' \ar[r] \ar[d] & 0 \\
0 \ar[r] &
M'[f^\infty] \ar[r] &
M' \ar[r] &
M'/M'[f^\infty] \ar[r] & 0,
}
$$
the third vertical arrow is an isomorphism. Since the rows are exact and the
first vertical arrow is an isomorphism by
Lemma \ref{lemma-torsion-completion} and $M[f^\infty] = M'[f^\infty]$,
the five lemma implies that $M \otimes_R R' \to M'$ is an isomorphism.

\medskip\noindent
The above shows that $\text{Can}$ is essentially surjective and that
the functor $H^0$ maps into the category of glueable modules.
Due to the exactness of (\ref{equation-BL-cech-mod-re}) for glueable modules
we have $H^0 \circ \text{Can} = \text{id}$ on the
category of glueable modules. This implies $\text{Can}$ is fully
faithful by Lemma \ref{lemma-H0-Can-adjoint} combined with
Categories, Lemma \ref{categories-lemma-adjoint-fully-faithful}.
This finishes the proof.
\end{proof}